% GENERATED FILE. Edit canonical lessons, then run npm run generate:books.
% canonical-source-hash: fnv1a64:706318faafb98328
% canonical-lessons: PA-C255-dak, PA-C255-lifafa, PA-C255-mohar, PA-C255-kaghaz, PA-C255-chhapna

\chapter{Letters and Stamps}
\label{ch:pa-a2-letters-and-stamps}
\hlchaptermodality{255}

\begin{chapteropening}
\textbf{By the end of this chapter:} \emph{I can talk about sending a letter by post.}

Complete the last lesson of chapter 255: I can talk about sending a letter by post.
\end{chapteropening}

\section[\texorpdfstring{\pa{ਡਾਕ} (ḍāk)}{ਡਾਕ (ḍāk)}]{\pa{ਡਾਕ} (ḍāk) — the post; the mail}

\label{lesson:PA-C255-dak}



\begin{quote}
Before the new one: say the Punjabi for wrong, then the Punjabi for to rub out.
\end{quote}

\subsection*{You'll want to know: \pa{ਡਾਕ}}
\textbf{\pa{ਡਾਕ}} — \emph{ḍāk} — \textquotedblleft{}the post; the mail\textquotedblright{}.

\begin{quote}
\pa{ਡਾਕ} \pa{ਸਵੇਰੇ} \pa{ਦਸ} \pa{ਵਜੇ} \pa{ਆਉਂਦੀ} \pa{ਹੈ।}
\end{quote}

\begin{quote}
\emph{ḍāk savere das vaje āundī hai.}
\end{quote}

\begin{quote}
The post comes at ten in the morning.
\end{quote}

A feminine noun.

\begin{cousinweb}
\textbf{\pa{ਡਾਕ}} (\emph{ḍāk}) goes back to a reconstructed \emph{*ḍakka}, \textquotedblleft{}relay, post\textquotedblright{}; CDIAL gives no Punjabi form, but lists Hindi \emph{ḍāk}, \textquotedblleft{}relay, post\textquotedblright{} (CDIAL 5519).
\end{cousinweb}

\begin{grammarlens}[title={What you've built}]
One more word for reading notices and messages.
\end{grammarlens}

\subsection*{Your turn}
\begin{itemize}
\raggedright
  \item \emph{Say it:} \emph{ḍāk}
  \item \emph{Say it:} \emph{ḍāk}, once more
  \item \emph{Say it:} \emph{ḍāk savere das vaje āundī hai.} — The post comes at ten in the morning.
  \item \emph{Recall:} say the Punjabi for wrong, then the Punjabi for to rub out, then say \emph{ḍāk} again
\end{itemize}

\begin{tcolorbox}[breakable,colback=teal!4,colframe=teal!35!black,title={Before you move on}]
What does \pa{ਡਾਕ} mean? (\textquotedblleft{}The post; the mail\textquotedblright{}.) Say it once more.
\end{tcolorbox}
\section[\texorpdfstring{\pa{ਲਿਫ਼ਾਫ਼ਾ} (lifāfā)}{ਲਿਫ਼ਾਫ਼ਾ (lifāfā)}]{\pa{ਲਿਫ਼ਾਫ਼ਾ} (lifāfā) — an envelope}

\label{lesson:PA-C255-lifafa}



\begin{quote}
Before the new one: say the Punjabi for to rub out, then the Punjabi for the post.
\end{quote}

\subsection*{You'll want to know: \pa{ਲਿਫ਼ਾਫ਼ਾ}}
\textbf{\pa{ਲਿਫ਼ਾਫ਼ਾ}} — \emph{lifāfā} — \textquotedblleft{}an envelope\textquotedblright{}.

\begin{quote}
\pa{ਚਿੱਠੀ} \pa{ਲਿਫ਼ਾਫ਼ੇ} \pa{ਵਿੱਚ} \pa{ਰੱਖੋ।}
\end{quote}

\begin{quote}
\emph{chiṭṭhī lifāfe vich rakkho.}
\end{quote}

\begin{quote}
Put the letter in the envelope.
\end{quote}

A masculine noun.

\begin{cousinweb}
\textbf{\pa{ਲਿਫ਼ਾਫ਼ਾ}} (\emph{lifāfā}) is a loan, usually traced through Persian to Arabic \emph{lifāfa}, \textquotedblleft{}a wrapper\textquotedblright{}.
\end{cousinweb}

\begin{grammarlens}[title={What you've built}]
One more word for reading notices and messages.
\end{grammarlens}

\subsection*{Your turn}
\begin{itemize}
\raggedright
  \item \emph{Say it:} \emph{lifāfā}
  \item \emph{Say it:} \emph{lifāfā}, once more
  \item \emph{Say it:} \emph{chiṭṭhī lifāfe vich rakkho.} — Put the letter in the envelope.
  \item \emph{Recall:} say the Punjabi for to rub out, then the Punjabi for the post, then say \emph{lifāfā} again
  \item \emph{Write it:} copy \textbf{\pa{ਲਿਫ਼ਾਫ਼ਾ}} once, checking the whole word
\end{itemize}

\begin{tcolorbox}[breakable,colback=teal!4,colframe=teal!35!black,title={Before you move on}]
What does \pa{ਲਿਫ਼ਾਫ਼ਾ} mean? (\textquotedblleft{}An envelope\textquotedblright{}.) Say it once more.
\end{tcolorbox}
\section[\texorpdfstring{\pa{ਮੋਹਰ} (mohar)}{ਮੋਹਰ (mohar)}]{\pa{ਮੋਹਰ} (mohar) — an office stamp; a seal}

\label{lesson:PA-C255-mohar}



\begin{quote}
Before the new one: say the Punjabi for the post, then the Punjabi for an envelope.
\end{quote}

\subsection*{You'll want to know: \pa{ਮੋਹਰ}}
\textbf{\pa{ਮੋਹਰ}} — \emph{mohar} — \textquotedblleft{}an office stamp; a seal\textquotedblright{}.

\begin{quote}
\pa{ਅਰਜ਼ੀ} \pa{ਉੱਤੇ} \pa{ਦਫ਼ਤਰ} \pa{ਦੀ} \pa{ਮੋਹਰ} \pa{ਹੈ।}
\end{quote}

\begin{quote}
\emph{arzī utte daftar dī mohar hai.}
\end{quote}

\begin{quote}
The application has the office stamp on it.
\end{quote}

A feminine noun.

\begin{cousinweb}
\textbf{\pa{ਮੋਹਰ}} (\emph{mohar}) is a loan, usually traced to Persian \emph{muhr}, \textquotedblleft{}a seal\textquotedblright{}.
\end{cousinweb}

\begin{grammarlens}[title={What you've built}]
One more word for reading notices and messages.
\end{grammarlens}

\subsection*{Your turn}
\begin{itemize}
\raggedright
  \item \emph{Say it:} \emph{mohar}
  \item \emph{Say it:} \emph{mohar}, once more
  \item \emph{Say it:} \emph{arzī utte daftar dī mohar hai.} — The application has the office stamp on it.
  \item \emph{Recall:} say the Punjabi for the post, then the Punjabi for an envelope, then say \emph{mohar} again
  \item \emph{Write it:} copy \textbf{\pa{ਮੋਹਰ}} once, checking the whole word
\end{itemize}

\begin{tcolorbox}[breakable,colback=teal!4,colframe=teal!35!black,title={Before you move on}]
What does \pa{ਮੋਹਰ} mean? (\textquotedblleft{}An office stamp; a seal\textquotedblright{}.) Say it once more.
\end{tcolorbox}
\section[\texorpdfstring{\pa{ਕਾਗ਼ਜ਼} (kāghaz)}{ਕਾਗ਼ਜ਼ (kāghaz)}]{\pa{ਕਾਗ਼ਜ਼} (kāghaz) — paper}

\label{lesson:PA-C255-kaghaz}



\begin{quote}
Before the new one: say the Punjabi for an envelope, then the Punjabi for an office stamp.
\end{quote}

\subsection*{You'll want to know: \pa{ਕਾਗ਼ਜ਼}}
\textbf{\pa{ਕਾਗ਼ਜ਼}} — \emph{kāghaz} — \textquotedblleft{}paper\textquotedblright{}.

\begin{quote}
\pa{ਮੇਰੇ} \pa{ਕੋਲ} \pa{ਕਾਗ਼ਜ਼} \pa{ਅਤੇ} \pa{ਕਲਮ} \pa{ਹੈ।}
\end{quote}

\begin{quote}
\emph{mere kol kāghaz ate kalam hai.}
\end{quote}

\begin{quote}
I have paper and a pen.
\end{quote}

A masculine noun.

\begin{cousinweb}
\textbf{\pa{ਕਾਗ਼ਜ਼}} (\emph{kāghaz}) is a loan, usually traced to Persian \emph{kāghaz}, \textquotedblleft{}paper\textquotedblright{}.
\end{cousinweb}

\begin{grammarlens}[title={What you've built}]
One more word for reading notices and messages.
\end{grammarlens}

\subsection*{Your turn}
\begin{itemize}
\raggedright
  \item \emph{Say it:} \emph{kāghaz}
  \item \emph{Say it:} \emph{kāghaz}, once more
  \item \emph{Say it:} \emph{mere kol kāghaz ate kalam hai.} — I have paper and a pen.
  \item \emph{Recall:} say the Punjabi for an envelope, then the Punjabi for an office stamp, then say \emph{kāghaz} again
  \item \emph{Write it:} copy \textbf{\pa{ਕਾਗ਼ਜ਼}} once, checking the whole word
\end{itemize}

\begin{tcolorbox}[breakable,colback=teal!4,colframe=teal!35!black,title={Before you move on}]
What does \pa{ਕਾਗ਼ਜ਼} mean? (\textquotedblleft{}Paper\textquotedblright{}.) Say it once more.
\end{tcolorbox}
\section[\texorpdfstring{\pa{ਛਾਪਣਾ} (chhāpṇā)}{ਛਾਪਣਾ (chhāpṇā)}]{\pa{ਛਾਪਣਾ} (chhāpṇā) — to print}

\label{lesson:PA-C255-chhapna}



\begin{quote}
Before the new one: say the Punjabi for an office stamp, then the Punjabi for paper.
\end{quote}

\subsection*{You'll want to know: \pa{ਛਾਪਣਾ}}
\textbf{\pa{ਛਾਪਣਾ}} — \emph{chhāpṇā} — \textquotedblleft{}to print\textquotedblright{}.

\begin{quote}
\pa{ਦਫ਼ਤਰ} \pa{ਹਰ} \pa{ਹਫ਼ਤੇ} \pa{ਫ਼ਾਰਮ} \pa{ਛਾਪਦਾ} \pa{ਹੈ।}
\end{quote}

\begin{quote}
\emph{daftar har hafte fāram chhāpdā hai.}
\end{quote}

\begin{quote}
The office prints forms every week.
\end{quote}

Past (I …): \pa{ਮੈਂ} \pa{ਫ਼ਾਰਮ} \pa{ਛਾਪਿਆ}, \emph{maiṁ fāram chhāpiā}.

\begin{cousinweb}
\textbf{\pa{ਛਾਪਣਾ}} (\emph{chhāpṇā}) goes back to a reconstructed \emph{*chapp}, \textquotedblleft{}press, cover, hide\textquotedblright{}; CDIAL lists Punjabi \emph{chāpṇā}, \textquotedblleft{}to print\textquotedblright{}, under it, as a word taken from Hindi \emph{chāpnā} (CDIAL 4994).
\end{cousinweb}

\begin{grammarlens}[title={What you've built}]
One more word for reading notices and messages.
\end{grammarlens}

\subsection*{Your turn}
\begin{itemize}
\raggedright
  \item \emph{Say it:} \emph{chhāpṇā}
  \item \emph{Say it:} \emph{chhāpṇā}, once more
  \item \emph{Say it:} \emph{daftar har hafte fāram chhāpdā hai.} — The office prints forms every week.
  \item \emph{Recall:} say the Punjabi for an office stamp, then the Punjabi for paper, then say \emph{chhāpṇā} again
  \item \emph{Write it:} copy \textbf{\pa{ਛਾਪਣਾ}} once, checking the whole word
\end{itemize}

\begin{tcolorbox}[breakable,colback=teal!4,colframe=teal!35!black,title={Before you move on}]
What does \pa{ਛਾਪਣਾ} mean? (\textquotedblleft{}To print\textquotedblright{}.) Say it once more.
\end{tcolorbox}
